\documentclass[11pt,letterpaper]{article}
\usepackage{amsmath,amssymb,amsthm}
\usepackage[margin=1in]{geometry}
\usepackage[hidelinks]{hyperref}
\hypersetup{pdftitle={Squareclass growth in structured rational sumsets},pdfauthor={Pan Yicheng},pdfsubject={Deductions from known theorems with AI research and drafting assistance}}
\usepackage{enumitem}
\setlength{\parskip}{4pt}
\setlength{\parindent}{0pt}
\setlist{itemsep=2pt,topsep=4pt}
\newtheorem{theorem}{Theorem}[section]
\newtheorem{proposition}[theorem]{Proposition}
\newtheorem{corollary}[theorem]{Corollary}
\newtheorem{lemma}[theorem]{Lemma}
\theoremstyle{remark}\newtheorem{remark}[theorem]{Remark}
\newcommand{\Q}{\mathbb Q}
\newcommand{\F}{\mathbb F}
\newcommand{\C}{\mathbb C}
\newcommand{\cl}{\operatorname{cl}}
\newcommand{\rk}{\operatorname{rank}}
\newcommand{\supp}{\operatorname{supp}}
\title{\bfseries Squareclass growth in structured rational sumsets}
\author{Pan Yicheng\\[3pt]\small With AI research and drafting assistance from OpenAI}
\date{2 October 2026 \quad Draft for review}
\begin{document}
\maketitle
\begin{abstract}
We give two consequences of established theorems that restrict certain approaches to large rational square-sum grids. First, among translates $1+u$ of rational $S$-units, a fixed nonzero squareclass occurs at most $2^{16|S|+16}$ times. This follows from the Beukers--Schlickewei unit-equation bound. Second, a rational set of size $M$ contained densely in a proper generalized arithmetic progression of fixed dimension $d$ has at least $M^{2/(5d)}$ squareclasses, up to a logarithmic factor, in every translate. This follows by slicing into arithmetic progressions and applying Bombieri--Zannier. The Green--Ruzsa form of Freiman's theorem gives a corresponding obstruction for sets with bounded additive doubling. We state all hypotheses and prove the deductions. Their originality has not been established. They do not resolve Erd\H{o}s Problem 885.
\end{abstract}

\section{Definitions and scope}
For $t\in\Q^*=\Q\setminus\{0\}$, let $[t]$ denote its image in the group $\Q^*/\Q^{*2}$, viewed as a vector space over $\F_2$. Thus $[xy]=[x]+[y]$, and $[x]=[y]$ exactly when $x/y$ is a nonzero rational square. Signs are included: a negative rational is not a square in $\Q$.

For a finite rational set $T$, define
\[
 T^*=T\setminus\{0\},\qquad
 c(T)=|\{[t]:t\in T^*\}|,\qquad
 h(T)=\dim_{\F_2}\langle[t]:t\in T^*\rangle.
\]
We set $c(\varnothing)=h(\varnothing)=0$. For every such set,
\begin{equation}\label{eq:colors}
 c(T)\le 2^{h(T)}.
\end{equation}
Here $c(T)$ counts the classes present, whereas $h(T)$ measures the dimension of their span. Multiplication by a nonzero rational preserves $c(T)$; it need not preserve $h(T)$. In fact $|h(\lambda T)-h(T)|\le1$, since $[\lambda t]=[\lambda]+[t]$.

For finite sets $A,B\subset\Q$, write $A+B=\{a+b:a\in A,b\in B\}$. The associated grid has one row for each distinct element of $A$ and one column for each distinct element of $B$. We are interested in families with $|A|=|B|=M\to\infty$, all $a+b\ne0$, and
\begin{equation}\label{eq:target}
 h(A+B)=o(\log M).
\end{equation}
Section~\ref{sec:context} explains why such a family would give arbitrarily large rational square-sum grids and a positive answer to the all-size assertion in Erd\H{o}s Problem 885. The present note supplies no such family and no general obstruction for arbitrary rational sets.

All logarithms denoted $\log$ are natural; $\log_2$ denotes base two. Constants $C_B\ge1$ and $\beta>0$ below are absolute constants furnished by the Bombieri--Zannier theorem, enlarged once if necessary to cover short progressions.

\section{Shifted rational units}
A rational $S$-unit, for a finite set $S$ of rational primes, is a number
\[
 u=\pm\prod_{p\in S}p^{e_p},\qquad e_p\in\mathbb Z.
\]
No restriction is imposed on these exponents or on the sizes of the primes.

\begin{theorem}\label{thm:units}
Let $2\in S$, let $s=|S|$, and let $U$ be a nonempty finite set of rational $S$-units with $-1\notin U$. For every $q\in\Q^*$,
\begin{equation}\label{eq:mult}
 |\{u\in U:1+u\in q\Q^{*2}\}|\le 2^{16s+16}.
\end{equation}
Consequently,
\begin{equation}\label{eq:unitrank}
 c(1+U)\ge |U|\,2^{-16s-16},\qquad
 h(1+U)\ge\log_2|U|-16s-16.
\end{equation}
The bound is independent of $q$, including its prime support and height.
\end{theorem}

\begin{proof}
Fix $q$. For each solution $1+u=qv^2$, where $v\in\Q^*$, set
\[
 K=\Q(\sqrt q),\quad \alpha=v\sqrt q,\quad
 x=\frac{1+\alpha}{2},\quad y=\frac{1-\alpha}{2}.
\]
Then $x+y=1$ and $xy=-u/4\ne0$. Let $S_K$ be the set of finite primes of $K$ above the rational primes in $S$, and let $E$ be the group of elements of $K^*$ that are units at every finite prime outside $S_K$.

At any such prime, $u$ is a unit, so $1+u$ is integral. The equality $\alpha^2=1+u$ implies that $\alpha$ is integral there. Since $2\in S$, the numbers $x,y$ are integral there as well. Their product $-u/4$ is a unit, so each is a unit. Thus $x,y\in E$. This argument does not add primes dividing $q$ to $S$.

Suppose first that $q$ is not a rational square. The nontrivial automorphism $\sigma$ of the quadratic field $K$ sends $\alpha$ to $-\alpha$. Therefore
\[
 (x,y)\in G=\{(\varepsilon,\sigma(\varepsilon)):\varepsilon\in E\}
 \subset(\C^*)^2.
\]
The map $E\to G$ is injective. By the classical $S$-unit rank formula, recalled in \cite[Section 1.1]{FM},
\[
 \rk G=\rk E=r_1(K)+r_2(K)-1+|S_K|\le1+2s.
\]
Here the ordinary unit rank of a quadratic field is at most one, and at most two finite primes lie above each rational prime.

Beukers--Schlickewei \cite[Theorem 1.1]{BS} bounds the number of solutions of $x+y=1$ in a rank-$R$ finitely generated subgroup of $(\C^*)^2$ by $2^{8R+8}$. Their theorem in fact holds in its divisible closure. Applying it to $G$ gives at most $2^{16s+16}$ ordered pairs. Distinct values of $u$ give distinct pairs, since $u=-4xy$. Choosing one of the two signs of $v$ for each $u$ suffices.

If $q$ is a rational square, then $x,y$ are rational $S$-units. The group of rational $S$-units has rank $s$, so its product with itself has rank $2s$. The same theorem gives the smaller bound $2^{16s+8}$. This proves \eqref{eq:mult}. Counting $U$ by its classes and using \eqref{eq:colors} proves \eqref{eq:unitrank}.
\end{proof}

If $-1$ is allowed in $U$, remove it before applying the theorem: its translate is zero, which has no multiplicative squareclass. The logarithmic inequality then uses $\log_2|U\setminus\{-1\}|$ and is asserted only when this set is nonempty.

\subsection{An affine normalization for grids}
\begin{corollary}\label{cor:anchor}
Let $A,B\subset\Q$ be finite, with $|A|\ge2$, $B\ne\varnothing$, and $a+b\ne0$ for every $(a,b)\in A\times B$. Fix any $a_0\in A$ and $b_0\in B$, and define
\begin{equation}\label{eq:anchor}
 U_{a_0,b_0}=\left\{\frac{a-a_0}{a_0+b_0}:a\in A\setminus\{a_0\}\right\}.
\end{equation}
Let $S$ contain $2$ and all primes appearing in the numerator or denominator of an element of $U_{a_0,b_0}$ in lowest terms. Then
\begin{equation}\label{eq:anchorbound}
 h(A+B)+16|S|+16\ge\log_2(|A|-1).
\end{equation}
The corresponding statement with $A$ and $B$ interchanged also holds.
\end{corollary}
\begin{proof}
The set $U_{a_0,b_0}$ has $|A|-1$ distinct nonzero elements. An element equals $-1$ only if $a+b_0=0$, which is excluded. For its element indexed by $a$,
\[
 1+u=\frac{a+b_0}{a_0+b_0},\qquad
 [1+u]=[a+b_0]+[a_0+b_0].
\]
Both classes on the right belong to the span of the grid classes. Thus $h(1+U_{a_0,b_0})\le h(A+B)$, and Theorem~\ref{thm:units} proves the claim.
\end{proof}

The normalized set in \eqref{eq:anchor} is unchanged under
\[
 A\mapsto\lambda A+t,\qquad B\mapsto\lambda B-t,
 \qquad \lambda\in\Q^*,\ t\in\Q,
\]
with the anchors transformed accordingly: both numerator and denominator acquire the factor $\lambda$. This invariance concerns $U$ and its prime support; the unnormalized grid span dimension may change under scaling.

In a family satisfying \eqref{eq:target}, the smallest permissible $S$ in every row normalization must satisfy
\[
 |S|\ge\left(\frac1{16}-o(1)\right)\log_2 M.
\]
The same is true for every column normalization. Thus sublogarithmic normalized prime support cannot provide such a family. Prime support is not the same as the abstract multiplicative rank of a group: one rational generator can involve arbitrarily many primes. No bound with that smaller rank substituted for $|S|$ is asserted here.

\section{Sets contained in proper progressions}
We first record the arithmetic-progression consequence needed for slicing.

\begin{lemma}\label{lem:AP}
Let $t_n=a+nv$ for $0\le n<L$, with $a,v\in\Q$, $v\ne0$, and integer $L\ge1$. For every $q\in\Q^*$,
\begin{equation}\label{eq:AP}
 |\{n:t_n\in q\Q^{*2}\}|\le C_B L^{3/5}(\log(2L))^{\beta}.
\end{equation}
The constants are independent of $a,v,q$.
\end{lemma}
\begin{proof}
Bombieri--Zannier \cite[Theorem 1]{BZ} gives this bound, with $\log L$ for $L\ge2$, for square entries of a nonconstant integer arithmetic progression. Divide the present progression by $q$ and choose a positive integer $D$ clearing the denominators of $a/q$ and $v/q$. Multiplication by $D^2$ gives an integer arithmetic progression. Every counted entry becomes a nonzero rational square that is an integer, hence an integer square. If necessary, reverse the progression and retain only its nonnegative terms; these form a contiguous progression of length at most $L$. The theorem applies to this progression, with the constant enlarged for lengths zero or one. This also covers negative $q$ and either sign of $v$.
\end{proof}

A rational generalized arithmetic progression is a set
\begin{equation}\label{eq:GAP}
 P=\left\{p_0+\sum_{i=1}^{d}n_iv_i:0\le n_i<L_i\right\},
 \qquad p_0,v_i\in\Q,\quad L_i\in\mathbb Z_{\ge2}.
\end{equation}
It is \emph{proper} if all the displayed representations are distinct. In that case its volume $V=\prod_iL_i$ equals $|P|$. Length-one directions can always be removed.

\begin{theorem}\label{thm:GAP}
Let $A\subset\Q$ have size $M\ge2$. Suppose $A\subset P$, where $P$ is the proper progression \eqref{eq:GAP}, $d\ge1$, and $V\le C_0M$ for some $C_0\ge1$. Put $L=\max_iL_i$. For any $b\in\Q$, let $z=1$ if $0\in A+b$ and $z=0$ otherwise. Then
\begin{equation}\label{eq:finiteGAP}
 c(A+b)\ge\frac{(M-z)L^{2/5}}{C_B V(\log(2L))^{\beta}}.
\end{equation}
In particular, define
\[
 \kappa(C_0)=\frac{1}{2C_BC_0(1+\log C_0/\log4)^{\beta}}>0.
\]
Then
\begin{align}
 c(A+b)&\ge\frac{\kappa(C_0) M^{2/(5d)}}{(\log(2M))^{\beta}},\label{eq:GAPcolors}\\
 h(A+b)&\ge\frac{2}{5d}\log_2M-\beta\log_2\log(2M)+\log_2\kappa(C_0).
 \label{eq:GAPrank}
\end{align}
These bounds permit a zero in the translate and are uniform in all rational steps, offsets, denominators, and shifts.
\end{theorem}
\begin{proof}
Fix a direction of length $L$, and fix the other $d-1$ coordinates. Properness partitions $P+b$ into exactly $V/L$ disjoint arithmetic progressions of length $L$. Their common step is nonzero, since $L\ge2$ and $P$ is proper. By Lemma~\ref{lem:AP}, any one class occurs at most
\[
 \frac VL C_B L^{3/5}(\log(2L))^{\beta}
   = C_B V L^{-2/5}(\log(2L))^{\beta}
\]
times in $P+b$, hence at most this many times in its subset $A+b$. There are $M-z$ nonzero elements to count. Division gives \eqref{eq:finiteGAP}. No density assumption on any individual slice of $A$ is needed.

Since $M\le V\le C_0M$,
\[
 L\ge V^{1/d}\ge M^{1/d},\quad L\le C_0M,\quad M-z\ge M/2.
\]
Moreover, for $M\ge2$,
\[
 \log(2L)\le\log(2M)+\log C_0
 \le\left(1+\frac{\log C_0}{\log4}\right)\log(2M).
\]
Substitution proves \eqref{eq:GAPcolors}. Taking base-two logarithms and using \eqref{eq:colors} proves \eqref{eq:GAPrank}.
\end{proof}

\subsection{Bounded additive doubling}
\begin{corollary}\label{cor:doubling}
For each fixed $K\ge1$ there are an integer $D(K)\ge1$ and a constant $C(K)$ such that, whenever $|A|=M\ge2$ and $|A+A|\le K M$, every $b\in\Q$ satisfies
\begin{equation}\label{eq:doubling}
 h(A+b)\ge\frac{2}{5D(K)}\log_2M
      -\beta\log_2\log(2M)-C(K).
\end{equation}
Consequently $h(A+b)\ge c_K\log M$ for all sufficiently large $M$, with $c_K>0$ depending only on $K$.
\end{corollary}
\begin{proof}
The Green--Ruzsa theorem \cite[Theorem 1.1]{GR} puts $A$ inside an actual proper coset progression in its ambient group, of dimension at most $D(K)$ and size at most $F(K)M$. Its stated estimates allow
\[
 D(K)=\max\{1,\lceil C K^4\log(K+2)\rceil\},\qquad
 F(K)=\exp\bigl(C K^4\log^2(K+2)\bigr)
\]
for a sufficiently large absolute $C$. The additive group $\Q$ has no nontrivial finite subgroup, so the finite subgroup part of the coset progression is trivial. We therefore obtain a proper rational progression as in Theorem~\ref{thm:GAP}, with $d\le D(K)$ and $C_0=F(K)$. Its actual rational values are retained; no change by an arbitrary Freiman isomorphism is being applied to squares. Equation~\eqref{eq:GAPrank} proves the result, taking $C(K)=-\log_2\kappa(F(K))$.
\end{proof}

If a balanced family satisfies \eqref{eq:target}, both
\begin{equation}\label{eq:diverge}
 \frac{|A+A|}{|A|}\longrightarrow\infty,
 \qquad \frac{|B+B|}{|B|}\longrightarrow\infty.
\end{equation}
Indeed, if the first ratio failed to tend to infinity, some infinite subsequence would have a fixed upper bound $K$. Choose any one element $b\in B$ for each grid on that subsequence. The span of its column classes is contained in the full grid span, so Corollary~\ref{cor:doubling} contradicts \eqref{eq:target}. Interchanging $A,B$ proves the other assertion.

\section{Connection with Erd\H{o}s Problem 885}\label{sec:context}
For a positive integer $N$, put
\[
 D(N)=\{|r-s|:r,s\text{ are positive integers and }rs=N\}.
\]
The all-size question asks whether, for every positive integer $k$, there are distinct positive integers $N_1,\ldots,N_k$ for which $\bigcap_{i=1}^kD(N_i)$ contains at least $k$ distinct numbers \cite{E885}. We include the elementary connections to identify exactly the scope of the preceding obstructions.

\begin{proposition}\label{prop:grid}
For a fixed $k$, the preceding assertion is equivalent to the existence of rational sets $A,B$ of sizes $k+1,k$, respectively, with all cross sums rational squares, allowing zero. Thus its truth for every $k$ is equivalent to the existence of arbitrarily large balanced rational square-sum grids whose entries are nonzero.
\end{proposition}
\begin{proof}
Given the integers and $k$ distinct common differences $d_j\ge0$, take
\[
 A=\{0,4N_1,\ldots,4N_k\},\qquad B=\{d_1^2,\ldots,d_k^2\}.
\]
For any factor pair $rs=N_i$ with $|r-s|=d_j$, one has $d_j^2+4N_i=(r+s)^2$. Every cross entry, including those in row zero, is a square.

Conversely, order the rows $a_0<a_1<\cdots<a_k$. For each $b_j$, choose nonnegative rational roots $r_j,s_{ij}$ with
\[
 r_j^2=a_0+b_j,\qquad s_{ij}^2=a_i+b_j.
\]
Let the positive integer $L$ clear all root denominators, and set
\[
 N_i=L^2(a_i-a_0),\qquad d_j=2Lr_j.
\]
The numbers $N_i$ are distinct positive integers, because each is the difference of the integer squares $(Ls_{ij})^2$ and $(Lr_j)^2$. Their positive integer factorizations
\[
 N_i=\bigl(L(s_{ij}-r_j)\bigr)\bigl(L(s_{ij}+r_j)\bigr)
\]
have differences $d_j$. These differences are distinct, since the distinct $b_j$ give distinct nonnegative $r_j$.

An arbitrarily large balanced square grid contains every required $(k+1)$-by-$k$ rectangle. Conversely the all-size assertion gives arbitrarily large balanced grids by restriction. To remove zeros, start with a $2m$-by-$2m$ square grid, select any $m$ rows, and discard the columns having a zero in these rows. Each row has at most one such column, so at least $m$ columns remain. This gives a zero-free $m$-by-$m$ grid.
\end{proof}

\begin{lemma}\label{lem:Ramsey}
Let $A,B$ each have $M$ distinct rational elements, with all cross sums nonzero. Suppose at most $r\ge1$ squareclasses occur. If
\begin{equation}\label{eq:Ramsey}
 M\ge\max\{2kr,\ k(2r)^k\},
\end{equation}
then some $k$-by-$k$ subrectangle has a single squareclass. Dividing both coordinate sets of this subrectangle by a representative of that class makes every entry a nonzero rational square.
\end{lemma}
\begin{proof}
Color each of the $M^2$ positions by its squareclass. Some color occurs at least $M^2/r$ times. Let $d_b$ be its degree at column $b$. Discrete convexity of $n\mapsto\binom nk$ on the nonnegative integers, with $\binom nk=0$ for $n<k$, gives
\[
 \sum_{b\in B}\binom{d_b}{k}\ge M\binom{\lfloor M/r\rfloor}{k}.
\]
Hence some $k$-element row set has at least
\[
 \frac{M\binom{\lfloor M/r\rfloor}{k}}{\binom Mk}
 \ge \frac{M}{(2r)^k}\ge k
\]
common columns of this color. For the middle inequality, every factor in the binomial ratio is at least $1/(2r)$ because $M/r\ge2k$. Selecting $k$ common columns proves the claim. If their class is $[q]$, each entry has the form $qv^2$, so $(a/q)+(b/q)=v^2$. This works for negative $q$ as well.
\end{proof}

Under \eqref{eq:target}, equation~\eqref{eq:colors} gives $r\le2^{h(A+B)}=M^{o(1)}$. For each fixed $k$, condition~\eqref{eq:Ramsey} eventually holds. For the original assertion at size $k$, apply Lemma~\ref{lem:Ramsey} at size $k+1$ and retain $k$ of its columns. Proposition~\ref{prop:grid} then gives the desired integers.

\section{What remains unresolved}
The preceding deductions require substantial structure. Theorem~\ref{thm:units} requires control of the prime support of a specific normalized coordinate set. Theorem~\ref{thm:GAP} requires containment with bounded relative volume in a proper progression of fixed dimension. Corollary~\ref{cor:doubling} supplies this containment when additive doubling is bounded.

Nothing here shows that arbitrary rational grids meet either structural condition. A possible family satisfying \eqref{eq:target} must have both the normalized prime-support growth from Corollary~\ref{cor:anchor} and the unbounded doubling in \eqref{eq:diverge}; these are necessary conditions, not a contradiction. Conversely, this note constructs no family in that remaining regime. In particular, the unrestricted question in Erd\H{o}s Problem 885 remains unresolved by this work.

\paragraph{Attribution and review status.}
The inputs are published results of Beukers and Schlickewei, Bombieri and Zannier, and Green and Ruzsa, together with the classical $S$-unit rank formula. The propositions here are presented as deductions from those results; no claim of originality or priority is made. AI assistance was used for literature retrieval, proof development, drafting, and consistency checks. This draft has not undergone peer review. Small exact-arithmetic checks of normalization identities, if supplied alongside it, are consistency tests rather than proofs of the cited theorems or of the unrestricted problem.

\begin{thebibliography}{9}
\bibitem{BS} F. Beukers and H. P. Schlickewei,
\emph{The equation $x+y=1$ in finitely generated groups},
Acta Arithmetica \textbf{78} (1996), 189--199, Theorem 1.1.
\href{https://doi.org/10.4064/aa-78-2-189-199}{doi:10.4064/aa-78-2-189-199}.
\href{https://webspace.science.uu.nl/~beuke106/s-units.pdf}{Author-hosted text}.

\bibitem{BZ} E. Bombieri and U. Zannier,
\emph{A Note on squares in arithmetic progressions, II},
Rendiconti Lincei. Matematica e Applicazioni, series 9,
\textbf{13} (2002), no. 2, 69--75, Theorem 1.
\href{https://www.bdim.eu/item?fmt=pdf&id=RLIN_2002_9_13_2_69_0}{Digitized primary publication}.

\bibitem{GR} B. Green and I. Z. Ruzsa,
\emph{Freiman's theorem in an arbitrary abelian group},
Journal of the London Mathematical Society (2) \textbf{75} (2007), 163--175, Theorem 1.1.
\href{https://doi.org/10.1112/jlms/jdl021}{doi:10.1112/jlms/jdl021}.
\href{https://arxiv.org/abs/math/0505198}{arXiv:math/0505198}.

\bibitem{FM} P. Fili and Z. Miner,
\emph{A generalization of Dirichlet's unit theorem},
\href{https://arxiv.org/abs/1210.7884}{arXiv:1210.7884},
Section 1.1, for the recalled classical $S$-unit rank theorem.

\bibitem{E885} T. Bloom, editor,
\emph{Erd\H{o}s Problem 885}, Erd\H{o}s Problems database.
\url{https://www.erdosproblems.com/885}.
\end{thebibliography}
\end{document}
